\section{Failure and null outcomes as persistent knowledge}

Failed, blocked, partial and undetermined executions are stored as addressable experience rather than discarded as execution noise. A failure episode distinguishes the observed mode, propagated effect, task/variation identity, competing causes, detection evidence, recovery actions, residuals and falsifiers. Negative outcomes remain in history after later success.

This makes failure a potential research input, but not an automatic method lesson. Let $e_j$ denote an immutable episode and $\phi(e_j)$ its observed failure signature. Repeated signatures can justify a recurrence hypothesis,
\[
\phi(e_1)\approx\phi(e_2)\approx\cdots\approx\phi(e_n),
\]
without justifying the causal statement that one mechanism caused every failure. Recurrence answers ``does something structurally similar persist?''; attribution answers ``which intervention should change the outcome?''

A valid development record therefore keeps failures, nulls and harmful interventions visible. Otherwise an adaptive system can manufacture an apparently improving history by retaining only successful challengers. Negative-history completeness is an assurance coordinate rather than an optional debugging convenience.

Where the record begins matters as much as what it holds. An execution record can expose a missing or failed coordinate of the method it was produced under; interpreting that coordinate as evidence \emph{about} the method is a separate step, and it is deliberately non-automatic, because a missing specification, a provider outage, a data defect and a genuine operator defect call for different responses.
